\documentclass[twoside, a4paper, 12pt]{article}
\usepackage{thesis-style}

\usepackage[backend=biber, style=trad-alpha]{biblatex}
\addbibresource{referencia.bib}

% Töltsd ki a saját szakdolgozatod adataival
\def\CIM{Az Erlang nyelv párhuzamos programozást támogató nyelvi elemeinek futtatható szemantikája}
\def\SZERZO{Palenik Mihály}
\def\VEDESEVE{2018}

\def\TANSZEK{Programozási Nyelvek és Fordítóprogramok Tanszék}
\def\TEMAVEZETO{Horpácsiné Kőszegi Judit}
\def\TEMAVEZETOBEOSZTAS{Tanársegéd, MSc}

\def\EXAMPLESFOLDER{source_examples}
\def\RULESFOLDER{tex/Szabalyok}

\def\T_COLOR{yellow}
\def\K_COLOR{green}
\def\DEFS_COLOR{blue}
\def\ENV_COLOR{red}
\def\PROCESSES_COLOR{orange}
\def\PROCESS_COLOR{magenta}
\def\PID_COLOR{BlueGreen}
\def\MAILBOX_COLOR{Plum}
\def\PROCESSFLAG_COLOR{Orchid}
\def\TERMINATEDPROCESS_COLOR{YellowGreen}
\def\PROCESSIDS_COLOR{RawSienna}
\def\PROCESSESCHECKMAILS_COLOR{Aquamarine}
\def\IDLEPROCESSES_COLOR{Apricot}
\def\PROCESSESDURINGEXIT_COLOR{CadetBlue}
\def\REGISTEREDPROCESSES_COLOR{Mahogany}
\def\LINKEDPROCESSES_COLOR{PineGreen}
\def\TIME_COLOR{Tan}

\title{\CIM}
\author{\SZERZO}
\date{\VEDESEVE}

\begin{document}
\pagestyle{empty}

% belső fedőlap
\begin{titlepage}

\begin{minipage}{0.40\linewidth}
\includegraphics[scale=0.3]{img/elte-cimer}
\end{minipage}
\begin{minipage}{0.50\linewidth}
\begin{center}
Eötvös Loránd Tudományegyetem \\
Informatikai Kar \\
\TANSZEK
\end{center}
\end{minipage}

\hrule
\vfill

\begin{center}
\Huge
\textbf{\CIM}
\normalsize
\end{center}

\vfill

\begin{minipage}[t]{0.5\linewidth}
\begin{flushleft}
\textbf{\TEMAVEZETO} \\
\TEMAVEZETOBEOSZTAS
\end{flushleft}
\end{minipage}
\begin{minipage}[t]{0.5\linewidth}
\begin{flushright}
\textbf{\SZERZO} \\
Programtervező Informatikus MSc
\end{flushright}
\end{minipage}

\vfill

\begin{center}
Budapest, \VEDESEVE
\end{center}

\end{titlepage}

\cleardoublepage

% a belső fedőlap utáni lap a témabejelentő

% tartalomjegyzék
\tableofcontents
\cleardoublepage

\pagestyle{plain}
\setcounter{page}{1}

% tartalom
% Ajánlott minden fő fejezetet külön fájlba írni, pl.:

%\section{Bevezetés}
Az idők folyamán egyre több nyelv jelent meg a szoftverfejlesztés területén, mely vagy egy adott problémakörre specializálódtak, vagy általános felhasználásra készültek. Amióta megjelentek a magasabb szintű programozási nyelvek, azóta szigorú nyelvtani szabályok alapján adható meg a program, mely hibák fordítási időben kiderülnek. Mostanra kidolgozott elméleti háttérrel rendelkeznek és egy újonnan létrejövő nyelvnél nem is kérdéses a formális szintaxis kidolgozása. Oka érthető, hiszen ezek alapján a szabályok alapján parszolják a kódot és építenek szintaxis fát, ami elengedhetetlen része egy fordítónak. Nagyobb probléma viszont a szemantikai hiba, melyet futásidőben veszünk észre. Gyakorlatban az idetartozó definíciókat informálisan szöveges leírással adják meg, annak ellenére, hogy a formális szemantika definiálására is kialakult már biztos elméleti háttér. Ennek kiküszöbölésére készülnek nagy mennyiségben tesztek, melyek próbálják kisebb-nagyobb sikerrel lefedni az összes lehetséges működést, és így a szemantikai hibákra fényt deríteni.

Másik megközelítés viszont a formális szemantikadefiníció készítése lehetne. Az informális definíciók nem pontosak, félreértésekhez vezethetnek. Mivel közvetlenül nem szokás formális szemantikai szabályokat használni fordítók írásakor, nehézkes is lenne, ezért nem készítenek. Ezenfelül egy-egy vezérlési szerkezetnek nagyon bonyolult, nehezen megfogalmazható szemantikadefiníciója van, ami még egy okot ad arra a nyelv készítőknek, hogy ne foglalkozzanak formális szemantika definiálásával. Sokszor még a nyelv fejlesztői között sincs egyetértés, hogy egy bonyolult de rövid programrészletnek mi az eredménye. Szükség van formális szemantikadefinícióra, és nem csak akadémiai körökben. Nélküle nem lehetne programhelyesség bizonyítást végezni, a kérdéses esetekben az informális definíciók nem tudnak kielégítő választ adni. Sőt mi több előfordulhatnak ellentmondások, vagy olyan szemantikai szabályok, ami a többi szabály miatt nem aktiválható.

A meglévő formális szemantikadefiníciókat viszont nehézkes alkalmazni közvetlenül az előbb felsoroltakra egységesen. Ezt próbálja meg kiküszöbölni a $\mathbb{K}$ keretrendszer, amivel operációs szemantikával megadhatjuk egy nyelvnek a formális szemantikáját, és a hozzátartozó többi eszköz segítségével közvetlenül tudunk a szemantika felhasználásával verifikálni.

A felsorolt problémák és a keretrendszer lehetőségei az, ami ösztönzött diplomamunkám elkészítésére. A következő fejezetekben az Erlang nyelv párhuzamos programozását támogató résznyelvének formális szemantikadefiníciója rajzolódik ki előttünk. Sikerült nagy részét lefednem a definícióval, ám kisebb hiányosságok maradtak, mint például a monitorozás, ám a folyamatok közti kapcsolat egymás felügyelésére megvalósult, amely jó kiindulást adhat ennek befejezésére.
%TODO
%	Kicsit lemásolni az előadásban mondottakat
%\include{tex/felhasznaloi}
%\include{tex/fejlesztoi}
%\include{tex/irodalom}


\section{Bevezetés}
dsfsd

\section{K keretrendszer}
dsfsdf

\section{Erlang nyelv}
Az \textit{Erlang} egy általános célra felhasználható funkcionális programozási nyelv. A Java nyelvhez hasonlóan egy szemétgyűjtővel rendelkező virtuális gépen fut a bájt-kódra lefordított program. Ez a fejezet egy rövid betekintést ad az Erlang nyelvhez a reference manual alapján (\cite{RefMan}). A lent található példákhoz nem szükséges fájlba mentett, modularizált kódot használnunk, közvetlenül az Erlang Shell segítségével interaktívan is ki tudjuk értékelni kifejezéseinket.

A következő alfejezetek az Erlang nyelv alapjait mutatja be, éppen annyit, amennyi elégséges lesz a szemantikadefiníciók megértéséhez. Mielőtt elkezdnénk az adattípusok áttekintését, ismerkedjünk meg a \textit{term} fogalmával. A term a kifejezések legegyszerűbb formája, tehát az adattípusok. Minden kifejezés végét ellenben a megszokott pontosvesszős megoldással, pont jelöli.

\subsection{Adattípusok}
Első körben érdemes az egyszerűbb adattípusokat áttekinteni, mielőtt megismernénk a különböző vezérlési szerkezeteket.

Két típusú szám literál létezik: \textit{integer} és \textit{float}. Alapértelmezetten decimális számrendszerben értendők, de a \textit{alap\#érték} formában 2-36-os számrendszerig ki tudunk fejezni számokat.

\lstinputlisting[language=Erlang]{\EXAMPLESFOLDER/numbers.erl}

A leggyakrabban használt literál az \textit{atom}, gyakorlatilag egy konstans. Az aposztróf nélküli atomoknak kis betűvel kell kezdődniük és alfanumerikus karaktereket, aláhúzást és @ jelet tartalmazhatnak. Az aposztróffal rendelkező atomok Latin-1 karakterkódolású karakterből állhatnak.

\lstinputlisting[language=Erlang]{\EXAMPLESFOLDER/atoms.erl}

A \textit{tuple} egy fix n-es, vagyis előre meghatározott számú termet tartalmazhat. A termeket kapcsos zárójelek közé, vesszővel elválasztva adjuk meg. Természetesen egymásba is ágyazhatóak.

\lstinputlisting[language=Erlang]{\EXAMPLESFOLDER/tuples.erl}

A \textit{lista} is a tuplehöz hasonlóan egy összetett típus ellenben változó hosszúságú. A listákat szögletes zárójelekkel adjuk meg, és szintén egymásba ágyazhatóak. Kétféleképpen lehet megadni Erlang nyelvben: a termek vesszővel elválasztott formában, vagy a fej elem, maradék lista formában.

\lstinputlisting[language=Erlang]{\EXAMPLESFOLDER/lists.erl}

Egy érdekes adattípus, ami manapság már sok imperatív nyelvben is megtalálható, a \textit{függvény objektum}. A \textit{fun} szóval tudunk létrehozni ilyen objektumot.

\lstinputlisting[language=Erlang]{\EXAMPLESFOLDER/fun.erl}

A példában egy függvény objektumot hoztunk létre, amely két paramétert vár X és Y, és ezeknek az összegével tér vissza. Láthatjuk is a példában, hogy a visszatérési értéke az objektumra mutató referencia. A nagy betűvel kezdődő szavak Erlangban a \textit{változók}, később még lesz róluk szó. A Sum változóban tárolt függvényt zárójeles formában meg is tudjuk hívni.

A \textit{pid} egy folyamatot azonosít három integer számmal.

\lstinputlisting[language=Erlang]{\EXAMPLESFOLDER/pid.erl}

A spawn függvény segítségével egy új folyamatot hozunk létre, és visszatérési értékként egy pidet ad vissza.

Még mielőtt folytatnák érdemes szót ejteni a \textit{sztring} és a \textit{logikai} literálokról. Erlangban mint külön adattípusok nem léteznek. A két logikai literál valójában két atom: a \textit{true} és a \textit{false}. A sztringek pedig listák. Idéző jelek között tudjuk megadni, de valójában az ASCII vagy unicode kódtáblában megfeleltetett számok listájaként van ábrázolva.

\subsection{Mintaillesztés}
Az előző részben már volt említés a változókról. A \textit{változó} egy kifejezés. Ha egy literál kötve van a változóhoz, akkor a visszatérési értéke az a literál. Minden változóhoz csak egyszer lehetséges literált kötni, ami történhet mintaillesztés során, vagy egyszerű értékadással (ami szintén egy mintaillesztés). A változók nagy betűvel vagy aláhúzással kezdődnek, és alfanumerikus karaktereket, aláhúzást vagy @ jelet tartalmazhatnak. Az aláhúzással kezdődő változóknak speciális jelentésük van. A fordító figyelmen kívül hagyja olyan értelemben, hogy nem generál warningokat a nem használt változók miatt.

\lstinputlisting[language=Erlang]{\EXAMPLESFOLDER/_variables.erl}

A példákban mintaillesztés látható. A \textit{minta} ugyan úgy épül fel mint egy term, ellenben tartalmazhat nem kötött változókat. \textit{Mintaillesztés} során pedig ezekhez a változókhoz kötünk értékeket. Ha a mintaillesztés sikertelen, mint a második példában, akkor \textit{badmatch} futásidejű hiba lép fel. A változókhoz csak egyszer köthető érték, és ez alól nem kivétel az aláhúzással kezdődőek sem, attól függetlenül, hogy a fordító figyelmen kívül hagyja.

Az anonymus változó csak egyetlen egy aláhúzásból áll. Akkor hasznos, ha kötelező változót megadnunk, de a benne lévő érték nem fontos.

\lstinputlisting[language=Erlang]{\EXAMPLESFOLDER/anonymusvariable.erl}

Az első példa jól mutatja, hogy gyakorlatilag figyelmen kívül hagyja az értékeket. A második példában egy lista fejelemére vagyunk kiváncsiak, amit mintaillesztéssel egyszerűen megkaphatunk.

Mintaillesztés nem csak az egyenlőség jellel lehetséges, hanem különböző vezérlési szerkezetekben is, mint például a \textit{receive} és a \textit{case}, amire később láthatunk példát.

\subsection {}

%TODO
%		Aztán kifejezések (listás bohóckodás stb
%		case, if, stb vezérlési szerkezetek
%		Pides része, mondani hogy itt nem térünk ki rá, majd a definícióknál lessz magyarázat
%		Megemlíteni, hogy van modulos felbontás, de nem szükséges ismerni, mert nincs rá még szemantika

%TODO változók



\subsection{Egy funkcionális nyelv}

\subsection{}


Milyen részei vannak
Hol használják főleg?
\section{Erlang nyelv konkurrens résznyelvének szemantikája K keretrendszerben}
vmi resz
\subsection{Erlang folyamatok}

    -Szemantika leírások
    -Itt beszélhetek tesztekről
    -Összehasonlítás más eddigi Erlang szemantikákkal
    
    Nagyvonalúan beszélni az erlang folyamatokról, a saját
szemantikadefiníció kifejtése. Külön részt szentelni a definíció megfogalmazását segítő, és
konzisztensen tartásáért felelős tesztekért.

\section{Verifikáció}
sdfsdfsd

\section{Utószó?}
fdsfsdfs


\newpage

\printbibliography[title = {Referencia}]


\end{document}
